\section{Interfaccia wan1 Aggiuntiva}\label{interfaccia-wan1-aggiuntiva}

Riprendendo l'architettura descritta in 6.1, per provare ad aggiungere
la connettività ad internet senza assegnare un'indirizzo gateway L3 al
bridge, e' possibile aggiungere ad ogni container duplicato, una
interfaccia \passthrough{\lstinline!wan1!} per l'inoltro verso
\passthrough{\lstinline!incusbr0!}.

Il problema di questa soluzione e' che altera le regole di routing dei
container, interferendo con il modo in cui il sistema nel container
instrada i pacchetti, eventualmente compromettendo la verifica
dell'exploit, oppure l'exploit stesso.

Inoltre e' una soluzione poco efficiente: richiede di duplicare
manualmente il numero di interfacce previste (tenendo a mente che ogni
coppia di interfacce richiede una coppia di veth per la comunicazione
con \passthrough{\lstinline!incusbr0!}). La soluzione sembra essere
abbastanza macchinosa da implementare per cui e' stata scartata.
